\section{Code and Implementation}
The core training pipeline lives in \texttt{Baselines/train\_ner\_baseline.py}. Launcher scripts in \texttt{Baselines/} select language, model and output directory, which keeps runs comparable and reduces the risk of configuration drift between experiments.

The shared script executes the following steps:
\begin{enumerate}
  \item read IOB2 data from disk,
  \item construct the label vocabulary,
  \item tokenize sentences with the chosen multilingual transformer tokenizer,
  \item align word-level labels to subword tokens,
  \item fine-tune the token-classification model,
  \item evaluate on the development set, and
  \item write predictions and metrics to the output directory.
\end{enumerate}

Label alignment deserves a brief note. Individual words can decompose into multiple subword tokens, so the implementation assigns the label to the first subword and masks subsequent ones in the loss — a common convention that keeps training stable and ensures that evaluation operates on token-level outputs consistent with the gold annotations.

The pipeline emits \texttt{dev\_metrics.json}, \texttt{dev\_predictions.iob2} and \texttt{test\_predictions.iob2} for every run. Having those artifacts made it straightforward to spot anomalies: sentence-count mismatches between prediction files and gold splits are immediately visible when you check them systematically (see Section~\ref{sec:analysis}).

The typological-distance computations are handled by a separate utility that operates on \texttt{lang2vec} WALS vectors. It is fully decoupled from model training and exists only to supply the structural evidence used to interpret observed transfer gaps.
